\documentclass[11pt,a4paper]{article}
\usepackage[margin=22mm]{geometry}
\usepackage{graphicx}
\usepackage[T1]{fontenc}
\usepackage{lmodern}
\usepackage{microtype}
\usepackage{booktabs}
\usepackage[colorlinks=true,linkcolor=black,urlcolor=black]{hyperref}
\setlength{\parskip}{0.6em}
\setlength{\parindent}{0pt}
\renewcommand{\arraystretch}{1.15}

\title{\vspace{-14mm}\textbf{Heuristics 2}\\[2mm]\large How the modular MECCG agent decides}
\author{}
\date{}

\begin{document}
\maketitle
\vspace{-14mm}

\section*{What the design is for}

Heuristics 1 scored moves with hand-tuned weights applied at the point of use.
The weights were not comparable across evaluators, so no number could be
checked against another and no explanation could be trusted. H2 exists to fix
one thing: \textbf{every quantity is in the same unit, and every hand-chosen
number has a name.}

That unit is \emph{tournament score differential} (TSD) --- the official scoring
gap --- converted into a win probability by a fitted model. A module says how
much a play changes TSD; the win-probability model says what that is worth in
this position. Anything not derived from card data, the player view or a
probability table must be a named field in \texttt{core/tunables.ts} and must
appear in the rationale tree of any decision that used it.

\section*{One decision, end to end}

\begin{center}
\includegraphics[width=\textwidth]{pipeline.pdf}
\end{center}

Three filters run before any module sees a candidate. \texttt{forwardActions}
drops what the engine has \emph{proved} undoes this phase's progress;
\texttt{cycleGuard} drops what has \emph{demonstrably} led back to this position,
which the engine cannot see because every lap mints fresh instance IDs; the plan
layer then re-ranks what survives.

Ownership is per decision, not per action. A module declares the action types it
owns and a context gate (\texttt{claims}); \texttt{pass} is owned by several
modules because it means something different in each place. A module that
declines any candidate on a decision withdraws from the whole decision, because
a partially scored list cannot be ranked honestly.

\textbf{There is no fallback.} H2 used to hand decisions it could not rank to
Heuristics 1, which was measured deciding 41.2\% of every real choice --- so
more than two decisions in five credited to H2 were H1's, and every module
improvement was measured on the fraction that reached the modules at all. A
ceiling made of another agent cannot be raised by working on this one, so it was
removed. When the ranking does not discriminate, H2 now plays one of the tied
best uniformly at random. \texttt{pass} wins a tie only when nothing else is
tied with it: a turn spent passing plays no resource, attempts no faction and
scores nothing, while the opponent's turn arrives regardless.

\clearpage
\section*{The three layers}

\begin{center}
\includegraphics[width=0.97\textwidth]{layers.pdf}
\end{center}

The layering is enforced by a test, not by convention:

\begin{itemize}\itemsep2pt
\item \textbf{modules} own action types and score candidates. They may
      \emph{not} import each other --- two modules that need the same
      valuation must share a service, so the number cannot drift apart.
\item \textbf{services} hold valuations more than one module needs. They may
      not import modules.
\item \textbf{core} holds units, the win-probability model, the constants and
      the rationale tree. It may not import services.
\end{itemize}

Which service each module actually reaches for --- the detail the arrows above
summarise:

\begin{center}\small
\begin{tabular}{ll}
\toprule
module & services it uses \\
\midrule
\texttt{travel} & beliefs, budget, defence, exposure, reach, strike/prowess \\
\texttt{combat} & character-value, standing, strike/\{prowess, sequence, model\} \\
\texttt{hazards} & attack-modifiers, beliefs, denial, exposure, hazard-plan, strike/* \\
\texttt{characters} & budget, character-value, defence, strike/prowess \\
\texttt{resources} & budget, character-value, defence, reach, strike/prowess \\
\texttt{factions} & budget, defence, reach, strike/prowess \\
\texttt{hand} & card-price, character-value, hazard-plan \\
\texttt{grants} & character-value, strike/prowess \\
\texttt{corruption} & character-value \\
\texttt{fetching} & card-price \\
\texttt{events} & event-value \\
\texttt{kill} & attack-value \\
\texttt{endgame}, \texttt{health} & --- (score from the view alone) \\
\bottomrule
\end{tabular}
\end{center}

Services depend on each other too, and that is where the sharing shows:
\texttt{card-price} is built from \texttt{budget}, \texttt{event-value},
\texttt{hazard-plan} and \texttt{standing}; \texttt{hazard-plan} from
\texttt{beliefs}, \texttt{denial}, \texttt{exposure} and the strike
enumeration; \texttt{defence} from \texttt{attack-modifiers},
\texttt{standing} and the same enumeration. Every one of those bottoms out in
\texttt{standing}, which is the only thing that converts marshalling points
into TSD.

\section*{What each part is worth}

Not every module is worth the same, and the difference is measured rather than
estimated. Flattening a group's opinion --- the candidates stay,
every one scores alike, so the tie rule picks among them uniformly --- and
re-running a 384-game gate gives what \emph{knowing which one to take} is worth:

\begin{center}
\begin{tabular}{lrr}
\toprule
opinion flattened & paired Elo (95\% CI) & cost \\
\midrule
--- (unmodified) & $-134$ [$-173$, $-97$] & --- \\
\texttt{play-hazard}, \texttt{place-on-guard} & $-166$ [$-208$, $-128$] & $-11$ \\
\texttt{assign-strike}, \texttt{resolve-strike}, \texttt{choose-strike-order} & $-202$ [$-247$, $-163$] & $-47$ \\
\texttt{plan-movement}, \texttt{enter-site}, \texttt{select-company} & $-249$ [$-299$, $-207$] & $-94$ \\
\bottomrule
\end{tabular}
\end{center}

\textbf{The hazard machinery is the most elaborate thing in H2 and it is inert.}
A bundle beam search resolves whole attacks state by state, \texttt{hazard-plan}
assigns every hazard in hand to a company, and creature prices are derived from
that plan --- replacing all of it with a coin flip costs 11 Elo, inside noise.
Movement carries 94 and had received almost no attention, because H2 already
agreed with human play about it often enough that it never rose up a
divergence-ranked list.

\section*{Where it stands, and what that means}

Against Heuristics 1 over 384 paired, side-swapped games: \textbf{31.6\%}, paired
Elo $-134$ [$-173$, $-97$]. Against human players the recorded corpus is
\textbf{107--0}, at a median 42 marshalling points to the AI's 2.

That second number is the one that matters, and it is not a margin. A side
losing 42--2 is not making slightly worse choices; it is never completing a
scoring sequence. The whole of H2's measurable opinion is worth about 140 Elo,
so an agent with no opinions at all would sit near $-300$ --- which bounds what
refining these modules can achieve, and says the remaining gap is in what H2
does not model at all rather than in what it models imperfectly.

\section*{How to look at it}

Every claim above came out of one of these, and none of them is a simulation of
taste --- each answers a different question, and mixing them up is how this
project has gone wrong before.

\begin{center}\small
\begin{tabular}{ll}
\toprule
instrument & the question it answers \\
\midrule
\texttt{explain} & why did H2 play \emph{that}, as a rationale tree with every tunable named \\
\texttt{coverage} & which action types have an owner at all \\
\texttt{human-compare} & does H2 choose what the human chose, per decision \\
\texttt{route-compare} & does the \emph{sequence} look like the human's \\
\texttt{scoring-loop}, \texttt{hand-flow} & where do marshalling points fail to arrive \\
\texttt{gate} & is this version stronger, over paired side-swapped seeds \\
\texttt{sweep} & what does this tunable do to the result \\
\bottomrule
\end{tabular}
\end{center}

Two cautions that were learned expensively.

\textbf{Agreement locates, it does not decide.} Corpus agreement found a draft
decided by coin flip, a placeholder price deciding the discard, and the 41\%
deferral to H1 --- all real defects, none visible from a win rate. But three
changes that \emph{raised} agreement cost 41, 42 and 84 Elo, while acting
on ties \emph{lowered} agreement nine points and gained 110. The human's answer
transfers only when it rests on something both players share: a rule binds both,
a strategy does not.

\textbf{Every corpus measurement is off-policy.} It replays positions a human
reached. H2 never plays those games --- in its own it reaches different, worse
positions, and its failures live there. That is why the two changes that paid
were defects visible \emph{regardless} of position: a decision made at random,
and an object never modelled at all.

\end{document}
